\documentclass[UTF8,a4paper,11pt]{ctexart}
\usepackage[margin=2.4cm]{geometry}
\usepackage{amsmath}
\usepackage{booktabs}
\usepackage{longtable}
\usepackage{array}
\usepackage{graphicx}
\usepackage{xcolor}
\usepackage{tikz}
\usetikzlibrary{arrows.meta,positioning,shapes.geometric,calc,fit,backgrounds}
\usepackage{pgfplots}
\pgfplotsset{compat=1.17}
\usepackage{listings}
\usepackage{tcolorbox}
\usepackage{hyperref}
\hypersetup{colorlinks=true,linkcolor=blue!55!black,urlcolor=teal!70!black,
  pdftitle={Aitachi-Hands 灵巧手全栈教材},pdfauthor={Aitachi}}
\usepackage{fancyhdr}
\pagestyle{fancy}\fancyhf{}
\fancyhead[L]{\small Aitachi-Hands 全栈教材}\fancyhead[R]{\small\thepage}
\renewcommand{\headrulewidth}{0.4pt}

\definecolor{keygreen}{RGB}{7,95,87}
\definecolor{soft}{RGB}{240,247,244}
\definecolor{warn}{RGB}{170,80,30}
\newtcolorbox{keybox}{colback=soft,colframe=keygreen,boxrule=0.8pt,arc=2pt,left=6pt,right=6pt,top=4pt,bottom=4pt}
\newtcolorbox{warnbox}{colback=warn!6,colframe=warn,boxrule=0.8pt,arc=2pt,left=6pt,right=6pt,top=4pt,bottom=4pt}

\lstset{basicstyle=\ttfamily\small,breaklines=true,frame=single,
  backgroundcolor=\color{gray!6},columns=fullflexible,
  keywordstyle=\color{keygreen},commentstyle=\color{gray}}

\newcommand{\code}[1]{\texttt{#1}}

\title{\Huge\textbf{Aitachi-Hands 灵巧手全栈教材}\\[6pt]
\Large 硬件设计 · 安装标定 · 软件对齐 · 驱动 · MuJoCo / Isaac Lab 仿真 · 强化学习双任务}
\author{Aitachi}
\date{2026 年 11 月 · v1.0}

\begin{document}
\maketitle
\thispagestyle{fancy}

\begin{keybox}
\textbf{读本定位}：本教材以 ORCA Hand v2 开源灵巧手为硬件本体，以飞特（Feetech）串行总线舵机为执行器，
贯穿\textbf{硬件设计 → 安装标定 → 软件对齐 → 驱动安装 → MuJoCo 仿真 → Isaac Lab 仿真 → 强化学习双任务对比实验}全链路。
所有命令、寄存器、参数表均取自本工程实测代码与日志（COM3 实机、17 轴标定、PPO v3 产线），
而非转抄手册。配套代码仓：\code{github.com/aitachi/Aitachi-Hand}。
\end{keybox}

\begin{keybox}
\textbf{命名对照（Aitachi-Hands 项目）}：本项目统一品牌为 \textbf{Aitachi-Hands}。文中
\code{AHand*}（AHandNative.exe、AHand-calibration.json、AHand/ 目录等）为实机工件的原生文件名，
一律对应 Aitachi-Hands 项目组件（如 AHandNative = Aitachi-Hands 原生控制台）；
\code{orca\_v2}/ORCA Hand v2 为上游开源硬件设计与官方资产目录名（CC BY 4.0，依许可保留署名）。
\end{keybox}

\tableofcontents
\newpage

\section{硬件设计}

\subsection{本体选型：ORCA Hand v2}
Aitachi-Hand 的机械本体采用开源灵巧手 \textbf{ORCA Hand v2}（ORCA Dexterity, Inc.，2025，arXiv:2504.04259），
许可证 CC BY 4.0（可商用）。选择它的三个理由：

\begin{enumerate}
  \item \textbf{全开源图纸}：CAD/STL/STEP/3MF 全量公开（393 文件 = 196 STL + 85 STEP + 15 3MF），
        每个 STL 配同名 STEP 源，可自行改设计；
  \item \textbf{成本可控}：论文口径整机 < 2000 CHF（$\approx$ ¥16,100），飞特舵机方案更低；
  \item \textbf{仿真资产齐}：官方提供 URDF + MJCF（v1/v2）+ 网格，MuJoCo/Isaac 开箱即用。
\end{enumerate}

\begin{keybox}
\textbf{核心参数速查}：17 DOF（16 手指 + 1 腕）· 滚动接触关节 · 0.4 mm 鱼线双腱驱动 ·
官方舵机位为 Feetech STS3215 或 Dynamixel 二选一（单总线 3 Mbps，3MF 内分两块电机板）·
PLA 结构件 + TPU 皮肤（或模具浇硅胶）· 06\_CNC 2 件需机加工 · 轴承/PTFE 管/M4 螺丝外购 · 单人装配 $\le$ 8 h。
\end{keybox}

\subsection{本机执行的差异化配置}
官方 BOM 给的是 STS3215/Dynamixel 位；Aitachi-Hand 实装机使用飞特 \textbf{HLS 系列舵机}（更新的固件线），
由 FD 上位机逐轴读出并归档（表~\ref{tab:servo-models}）。两种舵机指令帧格式同源（飞特 SCS/STS 总线协议族），
但寄存器布局与最大波特率不同，本教材第~\ref{sec:software}、\ref{app:servo} 章按 HLS 实测展开。

\begin{table}[h]\centering
\caption{实机舵机清单（FD.exe 读取，2026-11-01）}\label{tab:servo-models}
\begin{tabular}{ccccc}
\toprule
ID & 型号 & 型号编码 & 固件 & 行程原值 \\
\midrule
0 & HLS3930 & 4106 & 3.45 & 0--4095 \\
1--16 & HLS2915 & 6922 & 3.45 & 0--4095 \\
\bottomrule
\end{tabular}
\end{table}

\subsection{传动与结构要点}
\begin{itemize}
  \item \textbf{双腱驱动}：每关节两根 0.4 mm 鱼线（屈/伸各一），经 8 个绕线轮（05\_Spools）与 PTFE 导管走向舵机；
        装配增补文件 \code{ASSEMBLY\_ADDENDUM.md} 专门讲 PTFE 管长度、销钉与轴承装配，配有 17 张实拍图。
  \item \textbf{滚动接触关节}：指节间为滚动接触而非销轴，是 MJCF 关节范围的来源（见第~\ref{sec:software} 章映射表）。
  \item \textbf{皮肤}：09\_Skin 11 件 TPU 直打；或用 SiliconeMolds-1100.3mf 模具浇硅胶（触觉更佳）。
  \item \textbf{打印}：拓竹预排版成品文件 Prints-1100.3mf（含参数）；分件打印按 01\_Fingers(32)/02\_Carpals(12)/
        03\_Wrist(14)/04\_ForeArm(26) 取件。
\end{itemize}

\subsection{图纸与上游资源}

\begin{table}[h]\centering
\caption{上游资源索引}\label{tab:upstream}
\begin{tabular}{lll}
\toprule
资源 & 地址 & 用途 \\
\midrule
图纸仓 & github.com/orcahand/orcahand\_hardware & CAD/STL/STEP/3MF/装配增补 \\
描述仓 & github.com/orcahand/orcahand\_description & URDF/MJCF/网格（仿真入口） \\
主控仓 & github.com/orcahand/orca\_core（v0.12.0） & 官方 ROS 控制 \\
论文 & arxiv.org/abs/2504.04259 & 设计依据 \\
官网 & orcahand.com · orca.ethz.ch/github & 门户/镜像 \\
装配视频 & YouTube \code{TgIz7HiyaoU} & $\le$ 8 h 全程 \\
舵机手册 & doc.feetech.cn（FTServo e-manual） & 协议/寄存器/上位机 FD.exe \\
\bottomrule
\end{tabular}
\end{table}

\begin{warnbox}
\textbf{避坑}：codeload.github.com 限流时段可改走 api.github.com 或稍后重试；
国内下载官方仓建议整仓 zip 归档后本地解压（本机归档于 \code{AHand/1028/}，hardware 578 MB + description 745 MB）。
\end{warnbox}

\section{硬件安装与标定}

\subsection{装配流程（官方口径 $\le$ 8 小时）}
\begin{enumerate}
  \item \textbf{打印/机加工备料}：PLA 结构件（或直接 Prints-1100.3mf 成品板）、TPU 皮肤/硅胶浇注件、
        06\_CNC 2 件机加工、轴承与 M4 螺丝按 BOM；
  \item \textbf{指节装配}（01\_Fingers 32 件）：滚动接触关节对合 → 销钉固定 → 穿 0.4 mm 鱼线双腱；
  \item \textbf{腱路布设}：PTFE 管按 \code{ASSEMBLY\_ADDENDUM.md} 裁长，绕线轮预绕线，保持屈/伸两腱张力一致；
  \item \textbf{腕与掌座}（02\_Carpals 12 / 03\_Wrist 14 / 04\_ForeArm 26）；
  \item \textbf{舵机入位}：16 $\times$ HLS2915 进指/掌电机板位，HLS3930 进腕位；
  \item \textbf{电气}：单总线菊花链（注意终端匹配与共地），12 V 供电（实测工作区间 9.6--12.4 V）；
  \item \textbf{皮肤}：TPU 直打件热装或硅胶件浇注后套装。
\end{enumerate}

\subsection{舵机 ID 分配（关键工序）}
装柜前用 FD.exe 把 17 个舵机逐个改好 ID（1--16 手指 + 0 腕），\textbf{先单舵机接 USB 改号再上总线}，
避免同 ID 冲突。本机最终 ID--关节分配以实测位姿差分验证（表~\ref{tab:axis-map}）。

\subsection{逐轴行程标定}
标定目标：为每个轴记录 \textbf{min\_raw / center\_raw / max\_raw}，作为（a）安全限位、（b）角度换算基准、
（c）RL 目标合法性检查的依据。本机标定方法（FD.exe + 手动缓动，2026-11-10/11）：

\begin{enumerate}
  \item FD.exe 连接 COM3，回中全手（零位偏移寄存器不动）；
  \item 对轴 $i$ 缓慢推向两个机械极限，记录观测 \texttt{observed\_min/max\_raw}；
  \item 取内缩 50 raw 作安全限位：\texttt{min\_raw = observed\_min+50}, \texttt{max\_raw = observed\_max−50}；
  \center 值 = 观测区间中点（floor 取整）；相对坐标 \texttt{relative = raw − center\_raw}；
  \item 每轴 $\ge$ 7 个采样点入档（\code{AHand-calibration.json}，来源 CSV 在 \code{cli/sessions/}）。
\end{enumerate}

\begin{table}[h]\centering\small
\caption{17 轴标定总表（实测，2026-11）}\label{tab:axis-map}
\begin{tabular}{rllrrrrr}
\toprule
ID & 手指 & 机械关节（实测名） & min & center & max & rel 区间 & 采样 \\
\midrule
0 & 腕 & 手腕前后摇摆 & 891 & 1857 & 2823 & $\pm$966 & 11 \\
1 & 无名指 & 一级指根（左右） & 1802 & 2142 & 2483 & $\pm$340 & 11 \\
2 & 无名指 & 二级指根（前后） & 784 & 1521 & 2259 & $\pm$737 & 11 \\
3 & 无名指 & 中间关节 & 609 & 1766 & 2924 & $\pm$1157 & 9 \\
4 & 中指 & 指尖关节 & 1053 & 1656 & 2260 & $\pm$603 & 10 \\
5 & 中指 & 一级指关节（左右） & 1637 & 2016 & 2396 & $\pm$379 & 11 \\
6 & 食指 & 左右位移关节 & 1651 & 1926 & 2201 & $\pm$275 & 7 \\
7 & 食指 & 二级指根（前后） & 1735 & 2505 & 3276 & $\pm$770 & 8 \\
8 & 食指 & 指尖关节（前后） & 1153 & 2153 & 3153 & $\pm$1000 & 12 \\
9 & 小指 & 一级指根（左右摇摆） & 1666 & 2057 & 2448 & $\pm$391 & 10 \\
10 & 小指 & 二级指根（前后） & 1730 & 2477 & 3225 & $\pm$747 & 12 \\
11 & 小指 & 指尖关节 & 2189 & 2671 & 3153 & $\pm$482 & 12 \\
12 & 中指 & 二级指根（前后） & 1624 & 2467 & 3310 & $\pm$843 & 12 \\
13 & 拇指 & 二级指根（上下） & 2232 & 2289 & 2346 & $\pm$57 & 23 \\
14 & 拇指 & 三级关节（主屈轴） & 1304 & 1708 & 2112 & $\pm$404 & 21 \\
15 & 拇指 & 指尖关节 & 1125 & 1993 & 2861 & $\pm$868 & 22 \\
16 & 拇指 & 一级指根（左右） & 1899 & 2092 & 2286 & $\pm$193 & 156 \\
\bottomrule
\end{tabular}
\end{table}

\begin{warnbox}
\textbf{ID13 行程极窄（$\pm$57 raw）且采样 23 次}——该轴有既有机械阻力史；ID16 采样 156 次为复测轴。
14 号轴为拇指主屈轴（全行程约 1700 raw）。凡遇轴异常，先查腱张力与 PTFE 管走向，再查电流/温度。
\end{warnbox}

\subsection{验收基准}
\begin{itemize}
  \item 全轴 \textbf{卸力回读}：寄存器 40 写 0 后读回必须为 0（未确认=断电处理）；
  \item 三手势编码器到位容差 \textbf{$\pm$20 raw}；轨迹回放允许一次 $\le$40 raw 末端补正；
  \item 安全护栏：峰值电流 < 120 raw、温度 < 40\,\textdegree C、电压 9.6--12.4 V、停滞/静止高电流 500 ms 内切除。
\end{itemize}

\section{软件对齐设计}\label{sec:software}

\subsection{四层软件栈}
Aitachi-Hand 把"同一只手"表达在四个层面，任何动作定义自上而下穿透：

\begin{table}[h]\centering
\caption{软件栈分层}\label{tab:stack}
\begin{tabular}{clll}
\toprule
层 & 工具 & 角度表示 & 说明 \\
\midrule
上位机 & FD.exe（飞特官方） & raw 编码 0--4095 & 改 ID/零位/PID，手动缓动标定 \\
网页控制 & hand-control.py + HTML & raw 目标 + 档位 & 并行批次（$\le$3 轴）、电流/温度暂停 \\
原生 GUI & AHandNative.cpp（Win32/C++） & raw + 编译期轨迹库 & 三轴同步 SYNC WRITE、四重保护、参数调优 \\
仿真/RL & MuJoCo / Isaac Lab（本项目） & rad 关节角 & 训练与评估，导出真机目标 \\
\bottomrule
\end{tabular}
\end{table}

\subsection{关节名 ↔ 舵机 ID ↔ 仿真关节 映射表}
这是全项目最重要的对齐表（来源：2026-11-11 A$\to$B 位姿差分实测 + 摄像头矫正兜底）：

\begin{table}[h]\centering\small
\caption{MJCF 关节 ↔ 舵机 ID 映射（ORCA v2 右手）}\label{tab:mjcf-map}
\begin{tabular}{lllll}
\toprule
手指 & MJCF 关节 & 舵机 ID & 实测关节名 & 相对行程 \\
\midrule
腕 & \code{right\_wrist} & 0 & 手腕前后摇摆 & $\pm$966 \\
拇指 & \code{right\_t-abd} & 16 & 一级指根（左右） & $\pm$193 \\
拇指 & \code{right\_t-cmc} & 13 & 二级指根（上下） & $\pm$57 \\
拇指 & \code{right\_t-mcp} & 14 & 三级关节（主屈） & $\pm$404 \\
拇指 & \code{right\_t-pip} & 15 & 指尖关节 & $\pm$868 \\
食指 & \code{right\_i-abd} & 6 & 左右位移 & $\pm$275 \\
食指 & \code{right\_i-mcp} & 7 & 二级指根（前后） & $\pm$770 \\
食指 & \code{right\_i-pip} & 8 & 指尖（前后） & $\pm$1000 \\
中指 & \code{right\_m-abd} & 5 & 一级指关节（左右） & $\pm$379 \\
中指 & \code{right\_m-mcp} & 12 & 二级指根（前后） & $\pm$843 \\
中指 & \code{right\_m-pip} & 4 & 指尖关节 & $\pm$603 \\
无名指 & \code{right\_r-abd} & 1 & 一级指根（左右） & $\pm$340 \\
无名指 & \code{right\_r-mcp} & 2 & 二级指根（前后） & $\pm$737 \\
无名指 & \code{right\_r-pip} & 3 & 中间关节 & $\pm$1157 \\
小指 & \code{right\_p-abd} & 9 & 一级指根（左右） & $\pm$391 \\
小指 & \code{right\_p-mcp} & 10 & 二级指根（前后） & $\pm$747 \\
小指 & \code{right\_p-pip} & 11 & 指尖关节 & $\pm$482 \\
\bottomrule
\end{tabular}
\end{table}

\subsection{角度换算公式}
每轴线性映射（符号方向由 A/B 位姿差分确定）：
\begin{equation}
q_i \;=\; s_i \cdot \frac{\mathrm{raw}_i - \mathrm{center}_i}{\mathrm{max}_i - \mathrm{min}_i}\cdot \Delta q_i^{\max},
\qquad
\mathrm{raw}_i \;=\; \mathrm{center}_i + s_i \cdot \frac{q_i}{\Delta q_i^{\max}}\cdot (\mathrm{max}_i - \mathrm{min}_i)
\end{equation}
其中 $\Delta q_i^{\max}$ 取 MJCF 该关节限位范围 \code{jnt\_range}，$s_i \in \{+1,-1\}$ 为方向符号。
实操顺序：仿真侧用 MJCF 关节限位归一化；真机侧以三个锚点位姿（V/OK/OPEN 的 raw 目标表）做最小二乘定标，
再以摄像头指尖跟踪矫正残差（\code{deploy/camera\_align.py}）。

\subsection{三个锚点位姿（raw 目标，实测已验证）}
\begin{table}[h]\centering\small
\caption{锚点位姿 raw 目标（容差 $\pm$20；A=张开/回正，B=V 型）}\label{tab:poses}
\begin{tabular}{lrrrrrrrrrrrrrrrr}
\toprule
姿 & \multicolumn{16}{c}{舵机 ID 1--16 raw 目标} \\
势 & 1&2&3&4&5&6&7&8&9&10&11&12&13&14&15&16 \\
\midrule
A & 1899&1970&2000&1419&1966&2142&2308&1846&2437&2021&1963&2237&2181&2047&1906&1904 \\
B & 1899&1318&907&1419&1966&2142&2308&1846&2437&2774&2510&2237&2095&488&1715&1969 \\
\bottomrule
\end{tabular}
\end{table}

配套档位（三姿势共用，实测安全值）：加速度 200 · 速度 200 · 目标电流 200（软件电流暂停阈值）·
扭矩 200 · 温度暂停 50 · 到位容差 20 · 并行上限 3 轴。

\subsection{系统数据流总图}

\begin{center}
\begin{tikzpicture}[node distance=9mm,
  box/.style={draw=keygreen,fill=soft,rounded corners=2pt,align=center,inner sep=5pt,font=\small},
  hw/.style={draw=warn,fill=warn!6,rounded corners=2pt,align=center,inner sep=5pt,font=\small},
  arrow/.style={-{Stealth[length=2.5mm]},thick}]
  \node[box] (fd) {FD.exe 上位机\\标定/改号};
  \node[box,right=of fd] (cal) {AHand-calibration.json\\17 轴限位/中心};
  \node[box,right=of cal] (pose) {锚点位姿表\\A/B raw 目标};
  \node[box,below=of cal] (cpp) {AHandNative C++\\SYNC WRITE 轨迹库};
  \node[box,left=of cpp] (web) {hand-control.py\\网页三轴并行};
  \node[box,right=of cpp] (sim) {MuJoCo / Isaac Lab\\rad 关节角};
  \node[box,below=of sim] (rl) {RL 双任务\\PPO/TD3/SAC$\cdots$};
  \node[hw,below=of rl] (real) {实机 COM3 @1 Mbps\\17 舵机};
  \draw[arrow] (fd)--(cal); \draw[arrow] (cal)--(pose);
  \draw[arrow] (cal)--(web); \draw[arrow] (cal)--(cpp); \draw[arrow] (pose)--(cpp);
  \draw[arrow] (sim)--(rl); \draw[arrow] (rl) to[bend left=18] node[right,font=\scriptsize]{导出轨迹+分段回放} (real);
  \draw[arrow] (cpp)--(real); \draw[arrow] (web) to[bend right=12] (real);
  \draw[arrow,gray,dashed] (real) to[bend right=28] node[below,font=\scriptsize]{电流/温度/位置回读护栏} (cpp);
\end{tikzpicture}
\end{center}

\begin{keybox}
\textbf{对齐验收}：同一动作的实机 raw 曲线与仿真 rad 曲线（换算后）叠加对比；
先单轴叠加，再多指并行，再验回位误差与指尖几何（V 分指 $\ge$30 mm、OK 拇食指距 $\le$10 mm）。
\end{keybox}

\section{驱动安装与上位机}

\subsection{Windows 侧（实机控制线）}
\begin{enumerate}
  \item \textbf{USB 转串口驱动}：接入舵机总线转接板后设备管理器出现串口（CP210x/CH340 系列按转接板芯片装驱动），
       本机为 \textbf{COM3}，波特率 \textbf{1\,000\,000}（HLS 实测配置；官方 STS3215 线可到 3 Mbps）；
  \item \textbf{FD.exe 上位机}（飞特官方，网盘包 \code{FD.exe} 提取码 juds）：改 ID、调零位/PID、缓动标定、
        示波器式回读。\textbf{FD 与自研软件互斥占用 COM3}——运行 C++/Python 控制前必须先断开 FD；
  \item \textbf{AHandNative 原生 C++ 软件}：双击 \code{AHandNative.exe}（免安装，Win32 原生，不需要 Python）。
        按钮：连接检查 / V / OK / OPEN / 停止卸力 / 打开记录 + 三姿势四档位（加速度/速度/目标电流/扭矩，50--100\%）调优。
        命令行模式：\code{AHandNative.exe --selftest}（协议自检不开串口）、\code{--inspect}、\code{--run V|OK|OPEN}；
        从源码构建：\code{build.cmd}（MSVC，含 \code{hand\_cli.cpp} 与 \code{trajectory\_library.h}）；
  \item \textbf{Python 控制线}（网页三姿势）：\code{hand-control.py} + \code{hand-control.html}，
        依赖 \code{pyserial}；启动器 \code{启动手势控制.cmd} / \code{启动快速手势.cmd}。
\end{enumerate}

\begin{warnbox}
\textbf{安全铁律}：任何异常先"停止并卸力"（寄存器 40 $\to$ 0 并回读确认）；
若提示卸力未确认，立即断舵机电源。C++ 软件每次动作结束自动 16 轴卸力并读回。
\end{warnbox}

\subsection{C++ 软件安全护栏（实测档位）}
\begin{table}[h]\centering
\caption{AHandNative 四重保护（\code{guard()} 实测阈值）}\label{tab:guards}
\begin{tabular}{lll}
\toprule
护栏 & 阈值 & 动作 \\
\midrule
峰值电流 & $|$current$|$ $\ge$ 120 raw & 抛异常立即停 \\
温度 & $\ge$ 40\,\textdegree C & 抛异常立即停 \\
电压 & 出 9.6--12.4 V & 抛异常立即停 \\
静止高电流 & $\ge$77 raw 且 $|$速度$|$$\le$2 持续 500 ms & 切除（防堵转） \\
无进展 & 误差 $>$20 raw 超 1 s 无收窄 & 切除 \\
指令合法性 & 位置 50--4050/加速度 1--254/速度 1--600/电流 1--68/扭矩 1--1000 & 拒发 \\
\bottomrule
\end{tabular}
\end{table}

\subsection{WSL Ubuntu 侧（仿真与 RL 训练线）}
\begin{lstlisting}[language=bash]
# Ubuntu-22.04（WSL2），复用已验证 venv（torch 2.11.0+cu130 / SB3 2.9.2 / MuJoCo 3.3.7）
wsl -d Ubuntu-22.04
cd /mnt/c/Users/ASUS/Desktop/AHand/style1101
bash scripts/setup_ubuntu.sh          # 离线 wheels 优先安装
.venv/bin/python -m unittest discover -s tests -v
\end{lstlisting}

\begin{keybox}
共享文件 $\neq$ 共享 Python 环境：Windows 实机线（pyserial/原生 EXE）与 WSL 训练线（venv）各自独立，
互不安装对方依赖。Isaac Sim 用独立 venv \code{\$HOME/isaac\_sim/venv/bin/python}（见第 5.2 章）。
\end{keybox}

\subsection{快速手势面板}
\code{fast-gesture-panel.html} 为免安装轻量入口（一键 V/OK/OPEN，经本机 HTTP 控制服务），
调试现场优先用它，FD 与原生 GUI 留给标定与调优场景。

\section{仿真教程（一）：MuJoCo}

\subsection{安装与验证}
\begin{lstlisting}[language=bash]
pip install mujoco==3.3.7                    # 或 pip install -r requirements-lock.txt
python -c "import mujoco; print(mujoco.__version__)"
\end{lstlisting}
GPU 渲染需要 EGL/OSMesa 之一；纯数值训练与评估不依赖任何 OpenGL 后端（渲染失败不影响训练，只影响回放图）。

\subsection{模型资产}
\begin{itemize}
  \item 场景文件：\code{orca\_v2/scene\_right.xml}（含左右手 combined/left 单独场景）；
  \item 结构：\code{scene\_right.xml} $\to$ include \code{assets/scene.xml + options.xml + mjcf/orcahand\_right(.body).xml}，
        网格在 \code{assets/} 下；
  \item 原始来源：官方 \code{orcahand\_description} 仓 v2 目录（本机归档 \code{AHand/1028/}）；
  \item 官方预览：\code{python -m mujoco.viewer --mjcf=v2/scene\_combined.xml}。
\end{itemize}

\subsection{手势浏览器 Demo（实测配方）}
\begin{lstlisting}[language=bash]
cd AHand/mujoco
./.venv/Scripts/python.exe serve_demo.py --open        # Windows 原生
# 浏览器打开 http://127.0.0.1:8767 （不要直接双击 web/index.html）
\end{lstlisting}
按键：\code{1/V}=V 手势 · \code{2/O}=OK · \code{0}=张开 · \code{A}=自动交替 · 空格=暂停 · \code{R}=正视。
纯物理测试（无浏览器）：\code{gesture\_demo.py --test}；端口占用加 \code{--port 8768}。

\subsection{RL 环境加载（本项目 envs/）}
\begin{lstlisting}[language=python]
import gymnasium as gym, mujoco
import envs  # 注册 AitachiFree-v1 / AitachiSeq-v1
env = gym.make("AitachiSeq-v1")
obs, info = env.reset(seed=7001)
print(env.action_space.shape, env.observation_space.shape)   # (16,) (84,)
\end{lstlisting}
执行器设定与实机一致的伺服语义：位置执行器 \code{kp=12, kv=0.5}（腕 \code{kp=60, kv=3}），
物理步长 0.002 s、控制周期 0.02 s（每动作 10 物理步）、隐式快速积分器。

\subsection{常用命令速查}
\begin{table}[h]\centering
\caption{MuJoCo 线常用命令}\label{tab:mj-cmd}
\begin{tabular}{ll}
\toprule
目的 & 命令 \\
\midrule
原生查看器（模型） & \code{python -m mujoco.viewer --mjcf=scene\_right.xml} \\
浏览器手势 Demo & \code{serve\_demo.py --open}（127.0.0.1:8767） \\
无头物理自检 & \code{gesture\_demo.py --test} \\
RL 冒烟（free 任务） & \code{run\_matrix.py --task free --algo ppo --timesteps 2048} \\
渲染回放 GIF & \code{MUJOCO\_GL=egl demo.py V --model best\_model.zip} \\
\bottomrule
\end{tabular}
\end{table}

\begin{warnbox}
WSL 里 Vulkan 只枚举到 llvmpipe（软件设备）时，MuJoCo/Isaac 的原生窗口路径不可用；
用 EGL 离屏或浏览器服务渲染替代。回放失败必须如实记录，不得拿静态图冒充轨迹回放。
\end{warnbox}

\section{仿真教程（二）：Isaac Sim / Isaac Lab（Ubuntu·WSL）}

\subsection{数据流与路径}
\begin{center}
\begin{tikzpicture}[node distance=7mm,
  box/.style={draw=keygreen,fill=soft,rounded corners=2pt,align=center,inner sep=5pt,font=\small},
  arrow/.style={-{Stealth[length=2.5mm]},thick}]
  \node[box] (urdf) {URDF / STL\\原始资产};
  \node[box,right=of urdf] (usd) {物理 USD\\固定底座+17 关节};
  \node[box,right=of usd] (lab) {Isaac Lab\\CUDA PhysX + 测量 JSON};
  \node[box,right=of lab] (play) {FFplay / USDA\\轨迹回放导出};
  \draw[arrow] (urdf)--(usd); \draw[arrow] (usd)--(lab); \draw[arrow] (lab)--(play);
\end{tikzpicture}
\end{center}

\begin{itemize}
  \item 项目目录（WSL 视角）：\code{/mnt/c/Users/ASUS/Desktop/AHand}；
        原始模型 \code{1028/orcahand\_description/v2}，物理资产 \code{gpu-validation/usd}，
        圆棍场景 \code{gpu-validation/stick-grasp}；
  \item Isaac Python：\code{\$HOME/isaac\_sim/venv/bin/python}；Kit 配置：\code{ubuntu-validation/isaac\_physics.kit}；
  \item 硬件：RTX 3070 Laptop（CUDA 可用）；WSL Vulkan 仅 llvmpipe $\Rightarrow$ \textbf{走无界面求解 + 本地工具回放}。
\end{itemize}

\subsection{启动与依赖自检（整块可复制）}
\begin{lstlisting}[language=bash]
(
export AHAND_ROOT=/mnt/c/Users/ASUS/Desktop/AHand/gpu-validation
export AHAND_PY="$HOME/isaac_sim/venv/bin/python"
export OMNI_KIT_ACCEPT_EULA=YES
cd "$AHAND_ROOT" || exit 1
mkdir -p logs stick-grasp/runs || exit 1
/usr/lib/wsl/lib/nvidia-smi -L
test -x "$AHAND_PY" && echo "Isaac Python exists"
test -f usd/right/orca_right.usd && test -f usd/left/orca_left.usd && echo "Both USD assets exist"
)
\end{lstlisting}

\subsection{加载物理模型（左右手 GPU 动力学测试）}
\begin{lstlisting}[language=bash]
cd "$AHAND_ROOT"
bash ./run_isaac_orca_wsl.sh
\end{lstlisting}
成功标记：\code{ORCA\_ISAAC\_GPU\_DYNAMICS device=cuda:0} 与左右手各一条 \code{ORCA\_ISAAC\_OK}。
已验证结论：左右手各 17 关节在 CUDA PhysX 上通过 600 步测试，最大回位误差 0.01783 rad。
URDF 重新导入（按需）：URDF importer 启动后 2 分钟未产出即终止复测，实验请直接用已实测通过的 USD。

\subsection{圆棍交接主实验}
\begin{lstlisting}[language=bash]
cd "$AHAND_ROOT"
bash ./run_isaac_cylinder_handoff_wsl.sh      # 圆柱交接
bash ./run_isaac_stick_grasp_wsl.sh           # 圆棍抓取
tail -n 35 logs/handoff-latest.log            # 运行中另开终端看日志
\end{lstlisting}
当前实测：圆棍完成两次摇摆，抬升约 17.00 mm，摇摆最大漂移约 7.97 mm（超出验收 $\Rightarrow$ 记失败态），
保存 455 帧轨迹 JSON。该任务留作 Isaac 侧进阶练习。

\subsection{回放与导出}
\begin{lstlisting}[language=bash]
cd "$AHAND_ROOT"
OPENBLAS_NUM_THREADS=1 "$AHAND_PY" stick-grasp/render_ubuntu_handoff.py \
  --track stick-grasp/pad-sway-cushion.json \
  --output stick-grasp/Ubuntu-ORCA-Handoff.mp4 &&
ffplay -autoexit stick-grasp/Ubuntu-ORCA-Handoff.mp4     # 空格暂停/Q 退出

OPENBLAS_NUM_THREADS=1 "$AHAND_PY" stick-grasp/export_native_usd.py \
  --track stick-grasp/pad-sway-cushion.json \
  --output stick-grasp/Ubuntu-ORCA-Handoff.usda          # 导出 USDA
\end{lstlisting}

\subsection{实机 A/B 与仿真对照（网页教程专属能力）}
确认版教程 \code{AHand-IsaacSim-IsaacLab-Ubuntu-确认版.html} 的 \#hardware 节内嵌实机 A/B 历史串口采样：
按钮切换动作、下拉选轴，离线绘制该轴 raw 位置曲线（虚线=目标）、电流/负载/温度读数，
并有"位置目标 vs 旧标定区间"越界复核表——它是把仿真角度对齐到实机编码的日常工具。

\begin{keybox}
\textbf{口径提醒}：MuJoCo 的 V/OK 用模型关节目标做主实验；Isaac Lab 当前主线是圆棍交接。
两条仿真线共享同一 USD/MJCF 几何，但物理引擎与接触参数不同，跨引擎结论须分别验证。
\end{keybox}

\section{强化学习：双任务、算法族与对比实验}

\subsection{任务定义（与 \code{envs/} 一一对应）}

\subsubsection{任务一 \code{AitachiFree-v1}：五指全自由活动但不打架}
每回合在"自然位姿流形"上采样一个目标姿态：三锚点（OPEN/V/OK）的 Dirichlet 凸组合 + $\sigma{=}0.05$ rad 噪声，
裁剪进关节限位。要求全程自碰撞穿透 $\le$ 5 mm（失败线），\textbf{无碰撞（$\le$1 mm）保持 3 个控制周期}才算成功——
"不打架"被显式写进成功判据。回合预算 2.5 s，控制 16 指（腕固定）。

\subsubsection{任务二 \code{AitachiSeq-v1}：V $\to$ OK $\to$ OPEN 快速手势序列}
单回合内三相位接力：每相位须满足 v3 验收窗（位置误差 $\le$0.12 rad、速度 $\le$0.35 rad/s、
V 分指 $\ge$30 mm / OK 拇食距 $\le$10 mm、保持 3 周期）才切换目标；三相位全过=成功。
每相位预算 2.0 s（总 6.0 s）；观测 = 80 维关节块 + 3 维相位 one-hot + 1 维相位进度（84 维）。

\begin{center}
\begin{tikzpicture}[node distance=12mm,
  st/.style={draw=keygreen,fill=soft,rounded corners=3pt,align=center,inner sep=6pt,font=\small},
  arrow/.style={-{Stealth[length=3mm]},thick}]
  \node[st] (s0) {OPEN 起始\\$\sigma{=}0.02$ rad 扰动};
  \node[st,right=16mm of s0] (v) {相位 1：V\\验收窗+3 周期保持};
  \node[st,right=16mm of v] (ok) {相位 2：OK\\拇食距 $\le$10 mm};
  \node[st,right=16mm of ok] (op) {相位 3：OPEN\\回正五指};
  \node[st,below=14mm of ok,draw=warn,fill=warn!6] (fail) {失败：超时 / 穿透 $>$5 mm / 非有限态};
  \draw[arrow] (s0)--(v); \draw[arrow] (v)--node[above,font=\scriptsize]{$8{+}8\cdot$余时} (ok);
  \draw[arrow] (ok)--node[above,font=\scriptsize]{$8{+}8\cdot$余时} (op);
  \draw[arrow] (op)--node[above,font=\scriptsize]{成功 $+5$} ++(2.1,0);
  \draw[arrow,gray] (v) -- (fail); \draw[arrow,gray] (ok) -- (fail);
\end{tikzpicture}
\end{center}

\subsection{奖励函数（v3 数值验证权重）}
\begin{align*}
r =\;& -w_e\,\overline{e_{\mathrm{norm}}^2} - w_v\,\overline{v_{\mathrm{excess}}^2}
      - w_a\,\overline{(a_t - a_{t-1})^2} - w_f\,\overline{f_{\mathrm{norm}}^2}
      - w_c\,\mathrm{coll} - c_{\mathrm{time}} \\
    & + g_p\,(\bar e_{t-1} - \bar e_t) + b_{\mathrm{tol}}\mathbf{1}_{\mathrm{tol}}
      + r_{\mathrm{phase}} + \big(b_s + b_t(1 - t/T)\big)\mathbf{1}_{\mathrm{succ}}
      - p_f\mathbf{1}_{\mathrm{fail}}
\end{align*}
默认权重：$w_e{=}4,\,w_v{=}0.02$（速度只罚超 0.30 rad/s 部分）,\,$w_a{=}0.005,\,w_f{=}0.01,\,w_c{=}2$（穿透$-$0.5 mm 起算）,
$c_{\mathrm{time}}{=}0.02,\,g_p{=}2,\,b_{\mathrm{tol}}{=}0.8,\,b_s{=}b_t{=}10,\,p_f{=}2$；相位奖 $r_{\mathrm{phase}}{=}8{+}8\cdot$相位剩余比例。
\textbf{碰撞不提前终止}（持续惩罚）——历史上"早退自杀"曾被策略利用来逃避时间惩罚。

\begin{keybox}
\textbf{残差动作设计（v3 核心修复）}：动作不直接给目标，而是
\code{指令 = 归一化(目标先验 + 0.3×PPO 残差)}，且每周期限速 0.15 rad。
零策略即等价于可解的直接基线，PPO 在其上优化"更快+更稳+低压+低穿透"——
这也是本对比实验所有算法共享的公平起点。
\end{keybox}

\subsection{观测 / 动作 / 判据速查}
\begin{table}[h]\centering
\caption{两任务 MDP 规格对照}\label{tab:mdp}
\begin{tabular}{lll}
\toprule
 & 任务一 free & 任务二 seq \\
\midrule
观测 & 80 = 16$\times$5 块 & 84 = 80 + 相位 one-hot 3 + 进度 1 \\
动作 & 16 维 $\in[-1,1]$（残差） & 同左 \\
回合上限 & 2.5 s（125 步） & 6.0 s（300 步） \\
成功 & 误差+速度达标 且 穿透$\le$1 mm 保持 3 步 & 三相位全部验收 \\
关键指标 & 成功率 / 到位时间 / 最大穿透 & 成功率 / 序列总时 / 三相位分解 \\
\bottomrule
\end{tabular}
\end{table}

\subsection{16 篇论文先进性分类}
本工程从 70 篇方法总表（\code{PPO-方法总表.html}）中选取 16 篇，分四组进入"先进性阶梯"：

\begin{center}
\begin{tikzpicture}[
  lvl/.style={draw=keygreen,fill=soft,rounded corners=2pt,align=center,inner sep=4pt,font=\scriptsize},
  grp/.style={rounded corners=2pt,align=center,inner sep=5pt,font=\small\bfseries},
  arrow/.style={-{Stealth[length=2.5mm]},thick,dashed}]
  \node[grp,draw,fill=keygreen!12,text=keygreen] (g1) at (0,0) {A 基础与复现\\01 TRPO \textbullet{} 02 GAE \textbullet{} 03 PPO\\04 Impl-Matters \textbullet{} 05 OnPolicy-AC};
  \node[grp,draw,fill=blue!8,text=blue!60!black] (g2) at (8.3,1.5) {B PPO 算法改进\\06 TR-PPO-RB \textbullet{} 07 ESPO\\08 PPG \textbullet{} 25 ExO-PPO \textbullet{} 26 OPPO};
  \node[grp,draw,fill=orange!10,text=orange!70!black] (g3) at (8.3,-1.5) {C 环境奖励与探索\\09 PBRS \textbullet{} 10 Human-Pref \textbullet{} 11 RND\\12 Eureka \textbullet{} 13 TWTL \textbullet{} 27 Improve-RM};
  \node[grp,draw,fill=purple!8,text=purple!70!black] (g4) at (16,0) {D 落地参照\\31 TD3 \textbullet{} 32 SAC \textbullet{} 56 ANYmal\\55 RMA \textbullet{} 62 RL-100};
  \draw[arrow] (g1)--(g2); \draw[arrow] (g1)--(g3);
  \draw[arrow] (g2)--(g4); \draw[arrow] (g3)--(g4);
\end{tikzpicture}
\end{center}

\begin{longtable}{p{3.2cm}p{2.2cm}p{7.2cm}p{2.6cm}}
\caption{16 篇入选论文与本项目角色}\\
\toprule
论文 & 训练类型 & 核心机制（精读结论） & 本项目角色 \\
\midrule
\endfirsthead
\toprule
论文 & 训练类型 & 核心机制 & 角色 \\
\midrule
\endhead
01 TRPO & 基础与复现 & 旧策略轨迹上估计优势；样本平均 KL 约束 + 共轭梯度 + 回溯线搜索 & \textbf{基线 trpo（实跑，CG+线搜）} \\
02 GAE & 基础与复现 & $\delta_t$ 按 $(\gamma\lambda)^l$ 衰减求和；$\lambda{\in}[0.92,0.98]$ 最快 & 全体 PPO 优势估计 \\
03 PPO & 基础与复现 & $\min(\rho\hat A,\mathrm{clip}(\rho)\hat A)$；$\varepsilon{=}0.2$ 得 0.82 归一分 & \textbf{基线（实跑）} \\
— PPO-KL & PPO 变体 & 自适应 KL 惩罚 $\rho\hat A-\beta\,\mathrm{KL}$；$\beta$ 随 KL 出 $[0.5,1.5]\times$目标 倍增/减半 & \textbf{变体 ppo\_kl（实跑）} \\
57 ACER & 连续控制 & 截断重要性比率 + Retrace + 回放纠偏；SB3 实现仅离散动作 & 分类引用（离散对照，见跨研究表） \\
04 Impl-Matters & 基础与复现 & 16 组实现开关：初始化/奖励尺度/归一化贡献 17--21\% & 公平实验协议来源 \\
05 OnPolicy-AC & 基础与复现 & 25 万 agent 扫 50+ 设计选择；裁剪是强基线 & 公平实验协议来源 \\
06 TR-PPO-RB & PPO 改进 & 越界反向斜率 rollback；逐状态 KL 触发 & \textbf{变体 ppo\_rb（实跑）} \\
07 ESPO & PPO 改进 & 优化未裁剪目标 + 漂移阈值 0.25 早停（优于 KL 0.01/0.05） & \textbf{变体 ppo\_espo（实跑）} \\
08 PPG & PPO 改进 & 策略/价值相位分离 + 辅助阶段价值蒸馏 + KL 锚 & \textbf{变体 ppo\_pg（实跑）} \\
25 ExO-PPO & PPO 改进 & 扩展离策略下界；越界梯度随偏移平滑压缩不归零 & \textbf{变体 ppo\_exo（实跑）} \\
26 OPPO & PPO 改进 & 生成/打分重叠 + 超额调度；多节点步延迟 4.49$\times$ & 系统组（定性引用，不实跑：无 LLM 服务环） \\
09 PBRS & 奖励与探索 & $+\gamma\Phi(s'){-}\Phi(s)$ 势能抵消保持策略不变 & \textbf{变体 ppo\_pbrs（实跑）} \\
10 Human-Pref & 奖励与探索 & 片段偏好训奖励模型集成；$\sim$700 对逼真奖励 & 分类引用（需人标） \\
11 RND & 奖励与探索 & 固定随机目标网 + 预测误差内在奖励；双价值头 & \textbf{变体 ppo\_rnd（实跑，单头简化）} \\
12 Eureka & 奖励与探索 & LLM 采样可执行奖励 + GPU 并行评价 + 迭代 & 分类引用（留作扩展位） \\
13 TWTL & 奖励与探索 & 时序逻辑塑形 + 离线策略混合 & 分类引用（任务二天然分段） \\
27 Improve-RM & 奖励与探索 & OOD 检测 + KL 邻域探索主动收集偏好 & 分类引用（需人标） \\
\bottomrule
\end{longtable}

\subsection{跨研究对照（RLP 基准）}
同源方法在姊妹仓 Reinforcement-Learning-Policy 的自有基准上另有三列评测（E1/E2/E3）：
PPG 0.669$\pm$0.02/0.681/1.000$\pm$0.00，PPO-KL 0.658/0.785/0.992，PPO-RB 0.655/0.669/0.995，
PPO-Clip 0.652/0.620/0.996（E2 种子散布 $\pm$0.17），ESPO 0.648/0.755/0.568，TR-PPO-RB 0.625/0.754/0.694，
TRPO E3=1.000 真值（代价为 CG+线搜开销）；ACER 为离散对照。两基准任务不同，数值不可直接横比，仅作方法族参照。

\subsection{负结果归因（7 跑追试）}
TRPO 三种子全败的机制已归因：默认探索尺度（$\log\sigma_0{=}0$，为 PPO 族 $-1.5$ 的 4.7 倍）
在 1 万步内摧毁可解残差先验，10k--500k 全检查点 0.000（预算假说否证）；匹配 init 后任务一
10k 处 0.80 但被 CG 更新侵蚀至 0.04，任务二无论 init 均结构性失败。RND 在
$\lambda{=}0.5/0.05$ 双任务四连零——方向冲突（新颖性对抗保持窗）而非尺度问题。
追加：$\lambda{=}0$ 对照组双任务 1.000（0.30s/1.14s）——wrapper 与后端无害，
剂量-反应 $\{0{:}1.0,\ 0.05{:}0,\ 0.5{:}0\}$ 锁定内在信号为唯一病因；
TRPO 匹配 init 三种子复现“救回-再侵蚀”（10k $\geq$0.80 $\to$ 150k $\leq$0.52）；
seq SAC 熵消融（$-16 \to -4$）仍 0.000——结构性失败，非熵尺度效应。
数据：\code{hands-rl/results\_attribution/}（含漂移曲线 checkpoint\_eval.json 与汇总 json）。

\subsection{三种基线算法一页纸}
\begin{table}[h]\centering\small
\caption{TD3 / SAC / PPO 要点（连续控制三族）}\label{tab:three}
\begin{tabular}{p{2.4cm}p{4.4cm}p{4.4cm}p{3.4cm}}
\toprule
 & TD3（2018） & SAC（2018） & PPO（2017） \\
\midrule
类别 & 离策略确定性 AC & 离策略最大熵随机 AC & 在策略随机 PG \\
核心 & 双 Q 取小防过估 + 目标策略平滑 + 延迟 actor & 熵正则 + 双 Q + 自动温度 & 裁剪重要性比 + GAE + 多轮小批更新 \\
探索 & 外加截断噪声 & 熵驱动内生探索 & 高斯策略采样 \\
样本效率 & 高（大回放） & 高 & 低（数据一次性） \\
本工程超参 & lr 1e-3, batch 256, 噪声 0.1, 延迟 2 & lr 3e-4, batch 256, 目标熵 $-16$ & lr 3e-4, 8$\times$256 步/批, 10 轮 \\
\bottomrule
\end{tabular}
\end{table}

\subsection{PPO 变体目标函数几何对比}

\begin{center}
\begin{tikzpicture}
\begin{axis}[width=13.5cm,height=6.4cm,xlabel={$\rho=\pi_\theta(a|s)/\pi_{old}(a|s)$},
  ylabel={每样本目标（$\hat A=1$ 归一）},legend style={font=\scriptsize,at={(0.02,0.98)},anchor=north west},
  domain=0.5:1.7,samples=200,grid=both,grid style={gray!18},xmin=0.5,xmax=1.7,ymin=-0.05,ymax=1.75,
  xtick={0.8,1.0,1.2,1.333,1.5}]
  \addplot[thick,keygreen]{min(x,1.2)};
  \addlegendentry{PPO-Clip $\varepsilon{=}0.2$（越界梯度=0）}
  \addplot[thick,warn,dashed]{x - 1.5*max(x-1.2,0)};
  \addlegendentry{TR-PPO-RB（越界反向斜率拉回）}
  \addplot[thick,blue,domain=0.5:1.7,samples=400] {x <= 1.2 ? x : 1.2 + 0.2*tanh((x-1.2)/0.2)};
  \addlegendentry{ExO 平滑压缩（梯度衰减不为 0）}
\end{axis}
\end{tikzpicture}
\end{center}

\subsection{对比实验设计（公平协议，源自 04/05 两篇方法论）}
\begin{itemize}
  \item \textbf{矩阵}：2 任务 $\times$ 11 算法（含 TRPO 与 PPO-KL） \{$\,$ppo, trpo, ppo\_kl, td3, sac, ppo\_rb, ppo\_espo, ppo\_ppg, ppo\_exo, ppo\_pbrs, ppo\_rnd$\,$\} $\times$ 3 种子；
  \item \textbf{预算}：每跑 150k 环境步；同一 MuJoCo 模型与目标文件（SHA-256 入档 \code{config.json}）；
        网络统一 256$\times$256（PPG 因需共享主干同样 256$\times$256 共享）；
  \item \textbf{评估}：训练后固定 25 个评估种子（7001--7025）确定性策略复测；成功率/到位或序列时间/穿透为硬指标；
  \item \textbf{训练曲线}：Monitor 逐集日志（\code{episodes.csv}）+ PPO 诊断（比率均值/最大值、clip 率、近似 KL、
        ESPO 停止轮、PPG 辅助损失，\code{ppo\_diagnostics.json}）；
  \item \textbf{实现披露}：TR-PPO-RB 的逐状态 KL 用 $|\log\rho|$ 代理（SB3 不存旧分布参数）；
        RND 为单价值头简化版；ExO 用 tanh 软饱和复现"平滑压缩"精神——三者均写入论文与 README 偏差声明。
\end{itemize}

\subsection{复现命令}
\begin{lstlisting}[language=bash]
cd Aitachi-Hand
bash scripts/smoke_all.sh                                   # 13 条路径冒烟
bash scripts/run_matrix_all.sh                              # 全矩阵 54 跑（断点续跑）
python scripts/aggregate.py --results results               # 汇总 summary.csv/json
python scripts/make_figures.py --results results --out figures   # 12 张对比图
\end{lstlisting}

\appendix

\section{附录 A：飞特舵机控制指令集（HLS 实测版）}\label{app:servo}
资料来源：官方手册入口 doc.feetech.cn（FTServo e-manual）；以下寄存器与帧格式以
\code{AHandNative.cpp} / \code{hand\_cli.cpp} 在 COM3 @1 Mbps 实测为准（HLS2915/3930，固件 3.45）。

\subsection{帧格式}
\begin{center}
\texttt{FF FF $\mid$ ID $\mid$ LEN $\mid$ 指令 $\mid$ 参数... $\mid$ 校验和}\qquad
校验和 $=\,\sim\!\sum(\text{ID..末参数})$ 低 8 位
\end{center}
\begin{table}[h]\centering
\caption{指令集}\label{tab:instr}
\begin{tabular}{cll}
\toprule
码 & 名称 & 说明 \\
\midrule
0x01 & PING & 探活，返回状态包 \\
0x02 & READ & 读寄存器：参数 = 起始地址 + 字节数 \\
0x03 & WRITE & 写寄存器：参数 = 起始地址 + 数据 \\
0x04 & REG WRITE & 缓存写，等 ACTION 触发 \\
0x05 & ACTION & 触发缓存写 \\
0x06 & RESET & 恢复出厂（慎用，会清 ID/零位） \\
0x83 & SYNC WRITE & 多轴同步写（本机三轴一批，广播 0xFE） \\
\bottomrule
\end{tabular}
\end{table}

\subsection{关键寄存器（HLS 实测）}
\begin{table}[h]\centering
\caption{寄存器表（地址十进制）}\label{tab:regs}
\begin{tabular}{rll}
\toprule
地址 & 名称 & 读写要点 \\
\midrule
3 (word) & 型号编码 & 6922=HLS2915；4106=HLS3930；preflight 校验用 \\
9/11 (word) & 最小/最大位置限位 & 写入目标越界即拒发 \\
30 & 工作模式/分辨率 & 实测应为 1（0--4095 全程） \\
33 & — & 实测应为 0 \\
40 & 转矩使能 & 1=上电保持；0=卸力；\textbf{卸力必须回读确认} \\
41 (9B) & 运动指令 & 加速度(1B) + 位置(2) + 电流(2) + 速度(2) + 扭矩限制(2)，小端 \\
56 (15B) & 反馈区 & [0:2]位置 [2:4]速度 [6]电压 [7]温度 [13:15]电流 \\
\bottomrule
\end{tabular}
\end{table}
位置/速度/电流为\textbf{符号-数值 15 位编码}（最高位符号，低 15 位幅值）。

\subsection{三轴同步写示例（0x83，本机实测配方）}
\begin{lstlisting}
FF FF FE 22 83 29 09          # 广播ID=FE, LEN=4+10*3, 起始寄存器41, 每轴9字节
  <ID1> <accel> <pos_L pos_H> <cur_L cur_H> <spd_L spd_H> <lim_L lim_H>
  <ID2> ...                    # 每批最多 3 轴；ID 不得重复
  <ID3> ...
  <checksum = ~sum(ID..lim_H)>
\end{lstlisting}
合法性窗口（越界拒发）：位置 50--4050 · 加速度 1--254 · 速度 1--600 · 电流 1--68 · 扭矩限制 1--1000。
发送后逐轴 READ 寄存器 41 回读比对，再轮询 56 反馈直到连续 3 次 $|$位置$-$目标$|$$\le$20 且 $|$速度$|$$\le$2。

\subsection{FD.exe 上位机操作要点}
改 ID（单舵机接 USB）→ 读参数页（型号/固件/零位偏移）→ 缓动采集 \code{min/center/max} →
导出 CSV。注意 FD 与自研软件互斥 COM3；FD 的"零位偏移"写 EEPROM，
本项目标定只读不写，换算在软件侧完成。

\subsection{常用低层调用（\code{hand\_cli} 语义）}
\begin{lstlisting}[language=bash]
./AHandNative.exe --inspect          # 巡检：连 17 轴，识别当前姿态，不发运动
./AHandNative.exe --run V            # 执行 OPEN->V 轨迹（结束自动卸力回读）
./AHandNative.exe --selftest         # 协议/调优/护栏自检（不开串口）
python align_real.py --status        # 只读巡检（部署包）
python align_real.py OK --dry-run    # 算目标不下发
python align_real.py --probe 14 +60  # 单轴微动探针（步长上限 60 raw）
\end{lstlisting}

\section{附录 B：超参数总表}
\begin{table}[h]\centering\small
\caption{训练超参（与 \code{algorithms/baselines.py} 同步）}\label{tab:hp}
\begin{tabular}{lp{5.6cm}p{6.4cm}}
\toprule
族 & 参数 & 值 \\
\midrule
PPO 共用 & lr/$\gamma$/$\lambda$/裁剪/熵/值系数/梯度范数 & 3e-4 / 0.99 / 0.95 / 0.2 / 5e-4 / 0.5 / 0.5 \\
 & 采样/批量/轮数/网络 & 8 环境 $\times$ 256 步 / 512 / 10 / [256,256]，$\log\sigma_0{=}-1.5$ \\
TR-PPO-RB & rollback 系数 $c$ / KL 代理触发 & $(1{+}\varepsilon)/(1{-}\varepsilon){=}1.5$ / $|\log\rho|{>}0.5$ \\
ESPO & 漂移阈值 / KL 阈值 & 0.25 / 0.05（epoch 间早停） \\
PPG & 策略轮/价值轮/辅助轮/KL$\beta$/间隔 & 3 / 6 / 6 / 1.0 / 每 9 次 rollout \\
ExO & 软饱和 & $\varepsilon\tanh((\log\rho{-}\varepsilon)/\varepsilon)$ \\
PBRS & 势函数/尺度 & $\Phi{=}-2\overline{|e_{norm}|}$，$F{=}\gamma\Phi'{-}\Phi$ \\
RND & 隐层/嵌入/lr/$\lambda$ & 256 / 64 / 1e-4 / 0.5（单价值头简化） \\
TD3 & lr/缓冲/起始/batch/$\tau$/噪声/延迟 & 1e-3 / 300k / 5k / 256 / 5e-3 / $\sigma{=}0.1$ / 2 \\
SAC & lr/缓冲/起始/batch/$\tau$ & 3e-4 / 300k / 5k / 256 / 5e-3（目标熵 $-16$） \\
\bottomrule
\end{tabular}
\end{table}

\section{附录 C：一键复现}
\begin{lstlisting}[language=bash]
git clone https://github.com/aitachi/Aitachi-Hand && cd Aitachi-Hand
python scripts/fetch_assets.py            # 拉官方 orcahand_description 模型(111MB, 不入仓)
bash scripts/smoke_all.sh                 # 13 路径冒烟(~4 min)
bash scripts/run_matrix_all.sh            # 54 跑全矩阵(约 6-8 h, RTX3070 实测口径)
python scripts/aggregate.py               # results/summary.csv + summary.json
python scripts/make_figures.py            # figures/fig01..fig12 对比学术图
cd textbook && bash build_pdf.sh          # 本教材 PDF
cd paper && bash build_pdf.sh             # IEEE 论文 PDF
\end{lstlisting}

\section{附录 D：踩坑实录（按时间线）}
\begin{enumerate}
  \item \textbf{v2 零成功三重根因}：增量动作随机游走 / 采样噪声破坏保持判据 / 碰撞早退被利用。
        修复=绝对残差动作 + $\log\sigma_0{=}-1.5$ + 取消碰撞早退（改持续惩罚）；
  \item \textbf{\code{actuator\_forerange} 笔误}：MuJoCo 正确字段是 \code{actuator\_forcerange}；
  \item \textbf{WSL Vulkan=llvmpipe}：原生窗口渲染不可行 $\Rightarrow$ 无界面求解 + FFplay/浏览器回放；
  \item \textbf{SB3 2.9 API 变化}：on-policy 缓冲直出 torch 张量（勿 \code{np.clip}）、
        离策略 \code{net\_arch} 必须用 \code{\{pi,qf\}} 字典、优化器在 \code{policy.optimizer}；
  \item \textbf{WSL 后台任务}：\code{wsl\ \&\ nohup} 会被会话回收，需 Windows 侧常驻进程持有；
  \item \textbf{Git-Bash 路径转换}：\code{/mnt/c} 参数会被 MSYS 改写，包在 \code{bash -lc '...'} 内执行；
  \item \textbf{多版本交付}：每类产物只留一个最终版，更新=改单源重建原地覆盖。
\end{enumerate}


\end{document}
